\documentclass{kpi-lab}
\usepackage{biblatex}
\usepackage{enumitem}
\usepackage{longtable}
\addbibresource{sources.bib}

\title{Лабораторна робота \textnumero 2\\
  Реалізація алгоритмів генерації ключів гібридних криптосистем
}
\author{ФБ-51мн Подолянко Тимофій}
\date{}

\begin{document}
\maketitle
\setcounter{tocdepth}{1}
\tableofcontents
\bigskip

\textbf{Мета роботи} ---
дослідити принципи принципи формування криптографічно
стійких випадкових
послідовностей, генерації ключового матеріалу, тестування
та генерації простих чисел
для асиметричних криптосистем, а також основні вимог до
життєвого циклу
криптографічних ключів.

\chapter{Дослідження принципів генерації криптографічних
ключів}\label{chapter:rsa-keygen}

\begin{quote}\small\it
  Що таке ентропія у контексті генерації криптографічних
  ключів? Які джерела ентропії можуть використовувати
  сучасні програмно-апаратні системи?
\end{quote}

Ентропію в контексті генерації криптографічних ключів слід розуміти як
міру невпорядкованості, випадковості, мінливості закритої системи
\cite{Barker2020}.
Іншими словами, ентропія показує, наскільки значення певної
випадкової змінної, (в якості якої можна розглядати і ключ)
є непередбачуваним.

Більш загально, від доступної ентропії залежить рівень
стійкості захисту генератору випадкових бітів, або ключа та
пов'язаного із ним криптографічного алгоритму \cite{Barker2020}.

Джерела ентропії будуються на існуванні
фізичних або нефізичних процесів, вихід яких не можу бути
передбачений.
Такі процеси називають джерелом шуму, і вони лежать в
основі моделі джерела ентропії, яка розглядається в
\cite{Turan2018}.

В якості конкретних джерел шуму в \cite{Turan2018}
наводяться фізичні --- спеціалізоване обладнання,
та нефізичні --- системна інформація (вивід функцій API,
деякі дані з RAM, системний час),
або дані, введені людиною (наприклад, рухи миші).
Під \enquote{text}{системним часом} автори,, найбільш
ймовірно, мали на увазі
деякий лічильним, значення якого не прив'язується
безпосередньо до реального часу
і постійно змінюється.
Саме по собі це не є достатнім джерелом шуму.

\begin{quote}\small\it
  У чому полягає різниця між TRNG, PRNG та криптографічно
  стійким генератором псевдовипадкових чисел (CSPRNG)?
\end{quote}

TRNG (True Random Number Generator)
надає ентропію у вигляді випадкових чисел, отриманих через
спостереження випадкових фізичних процесів. \cite{Bartock2022}

PRNG (Pseudorandom Number Generator)
використовує початкове значення
та, можливо, додаткові входи для створення
послідовності бітів, яка виглядає як статистично незалежна
і незміщена, за допомогою певного алгоритма.

CSPRNG (Cryptographically Secure PRNG) підсилює вимоги до
PRNG через властивість, що вихід генератора непередбачуваний,
якщо початкове значення невідомо \cite{Barker2020a}.

\begin{quote}\small\it
  Яку роль виконує початкове значення (seed) генератора
  псевдовипадкових чисел? Які наслідки може мати можливість
  передбачити або відтворити seed?
\end{quote}

Початкове значення генератора визначає конкретну
створену псевдовипадкову послідовність.
Можливість передбачити чи відтворити seed дозволяє
перебачити також і створені псевдовипадкові значення.

\begin{quote}\small\it
  Чому звичайні генератори псевдовипадкових чисел,
  призначені для моделювання, ігор або статистичних задач,
  не повинні використовуватися для генерації криптографічних ключів?
\end{quote}

Саме тому, що звичайні генератори не мають цієї необхідної
для застосування у криптографії властивості.
Тому за декількома значеннями з псевдовипадкової послідовності
можна відновити початкове значення та передбачити наступні значення.
У випадку ключів, наприклад, період послідовності може виявитися
достатньо малим для відновлення всього ключа за його частиною.

\begin{quote}\small\it
  Яким чином сучасні операційні системи надають програмам
  криптографічно стійкі випадкові значення? Навести приклад
  відповідного API для однієї операційної системи або мови
  програмування.
\end{quote}

Наприклад, в ядрі Linux драйверам пристроїв надається інтерфейс для
використання різних подій в якості джерел шуму. Це, наприклад,
довільна інформація про пристрої, серійні номери, адреси;
час пошуку на блокових пристроях; виміри часу та
ідентифікатори джерел системних переривань.
Також окремо в якості джерела ентропії ядро може
використовувати апаратний генератор випадковості.

Доступ до псевдовипадкових даних здійснюється через  системний
виклик \texttt{getrandom()}, або спеціальні
\texttt{/dev/random}, \texttt{/dev/urandom}.
\texttt{/dev/urandom} надає вихід CSPRNG, в той час як
\texttt{/dev/random} надає доступ безпосередньо до джерела
ентропії системи.

\begin{quote}\small\it
  Чим принципово відрізняється генерація випадкового
  симетричного ключа від генерації пари ключів асиметричної
  криптосистеми?
\end{quote}

Випадковий симетричний ключ може бути
отриманий безпосередньо як вихід генератора випадкових бітів
заданої довжини, як описано у \cite{Barker2020}.

Ключи в асиметричній криптографії мають додаткові
математичні властивості, тобто неможливо застосувати
будь-який випадковий бітовий рядок в якості ключа.
Окрім властивостей, необхідних для математичної коректності операцій,
існують вимоги, зумовлені безпекою. Зокрема, для RSA у
\cite[][\pagename 405]{Knuth2021}
наводяться наступні: неподільність на 3 для початкових \(p\), \(q\);
\(p-1\) та \(q-1\)повинні мати принаймні один дуже великий дільник;
відношення \(\frac{p}{q}\) не повинно бути близьким до простого дробу;
інші вимоги.
Тому у \cite{Barker2020} для
генерації випадкового асиметричних ключів (пар ключів) вимагається
застосування спеціальних алгоритмів. В ці алгоритми одразу
закладено перевірку обраних доменних параметрів,
а згенеровані ключи гарантовано матимуть необхідні
математичні властивості.

В алгоритмах генерації зазвичай
можна виділити деякі криптографічні примітиви в якості складових.
Наприклад, для генерації ключів для криптографії на основі
факторизації цілих чисел необхідні примітиви генерації простих чисел,
тестування чисел на простоту.

\begin{quote}\small\it
  Чи гарантує використання алгоритму AES-256 рівень захисту
  у 256 біт, якщо ключ отримано з недостатньо випадкового
  або передбачуваного джерела? Обґрунтувати відповідь.
\end{quote}

Ні, не гарантує.
Як зазначено в \cite{Barker2020}, рівень стійкості захисту
ключи визначається, зокрема,
процесом його генерації. Таким чином, якщо процес генерації
ключа не мав рівень стійкості захисту
принаймні 256 бітів, то і ключ не матиме такого рівня
стійкості захисту.

\begin{quote}\small\it
  Які основні етапи життєвого циклу криптографічного ключа
  необхідно враховувати, крім безпосередньої генерації ключа?
\end{quote}

Етапи життєвого циклю ключа визначаються в тому числі
типом ключа встановленою політикою безпеки,
тобто можуть бути визначені по-різному.
Наприклад, у \cite{Barker2020a} пропонується поділ
життєвого циклу керування ключами на чотири етапи:
доопераційний, операційний, післяопераційний та \enquote{зруйнований}.

Серед станів виділяють, зокрема, на операційному етапі ---
активний, тимчасово припинений стан використання;
на післяопераційному етапі --- деактивований, скомпрометований стани.
Активний стан також передбачає різні можливі режими використання ---
для захисту та для обробки інформації.

Кожен етап може передбачати перебування ключа в одному з
багатьох можливих станів, та переходи між цими станами.

В для асиметричних ключів, життєві цикла публічного та
секретного ключів розглядаються окремо.

\chapter{Дослідження генерації простих чисел}

\begin{quote}\small\it
  Для чого при генерації ключів RSA необхідні великі прості числа?
\end{quote}

Рівень стійкості захисту в RSA ґрунтується на складності
задачі факторизації, а операції використовують властивості
простих чисел. Максимальний рівень стійкості захисту
ключів визначається довжиною модуля \cite{Barker2019}.
В свою чергу модуль визначається добутком двох секретних
простих чисел. Складність факторизації зумовлена розміром цих чисел.

\begin{quote}\small\it
  Чому на практиці для великих чисел використовуються
  ймовірнісні тести простоти?
\end{quote}

Відомі детерміновані тести, зокрема, \emph{AKS primality
testing}, занадто обчислювально складні
для перевірки за припустимий час чисел, що мають порядки,
достатні для забезпечення необхідного
рівня стійкості захисту.

\begin{quote}\small\it
  У чому полягає основна ідея тесту Ферма та чому
  позитивний результат цього тесту не гарантує простоту числа?
\end{quote}

Тест Ферма заснований на перевірці виконання малої теореми Ферма для
потенційного простого числа та випадкових цілих чисел.
Тест є ймовірнісним, тому кожна ітерація лише збільшує нашу
впевненість.
Ми не перевіряємо всі числа, а лише певну випадкову підмножину.
Також, мала теорема Ферма описує лише наслідок з простоти числа,
але не гарантує зворотній зв'язок, тобто деяке складене
число може задовольняти теорему.

\begin{quote}\small\it
  Яку проблему тесту Ферма частково вирішують більш
  досконалі ймовірнісні тести?
\end{quote}

Безліч складений чисел можуть задовольняти вираз \(b^{n-1}
\equiv 1 \mod n\).
Числа Кармайкла задовольняють його завжди, для всіх значень
\(b\), тобто їх тест Ферма не може виявити.
Інші ймовірнісні тести усувають цей недолік та можуть виявити
будь-які складені числа.

\begin{quote}\small\it
  Який принцип роботи тесту Міллера–Рабіна та яке значення
  має кількість його ітерацій?
\end{quote}

Тест Міллера-Рабіна перевіряє досліджуване число \(n\)
на порушення, за певною випадковою основою, малої теореми Ферма,
та опосередковано перевіряє наявність власних дільників в
\(n\) \cite{Rabin1980}.

Як відмічено у \cite[][\pagename 395]{Knuth2021}, ймовірність
виявлення складеного числа тестом Міллера-Рабіна
не залежить від значення цього числа,
та показано у \cite{Rabin1980}, що ймовірність невиявлення алгоритмом
складеного числа не перевищує \(\frac{1}{4}\).
Тому збільшення кількості ітерацій алгоритма до \(n\)
зменшує ймовірність помилки до \(\left(\frac{1}{4}\right)^n\).

\begin{quote}\small\it
  Чому дуже мала ненульова ймовірність помилки тесту
  простоти може бути прийнятною при практичній генерації
  криптографічних ключів?
\end{quote}

Відомі наразі детерміністичні тести простоти
мають надмірну обчислювальну складність для практичної реалізації,
або ґрунтуються на ґіпотезах.

У \cite{2023} пропонується розглядати
ймовірність \(p_{k,t}\) того, що число, яке пройшло \(t\)
ітерацій Міллера-Рабіна, є складеним (на відміну від
  ймовірності помилки тесту, під якою зазвичай розуміють
  ймовірність того, що число, яке пройшло \(t\) ітерацій
теста, є складеним).
Через оцінку цього значення підбирається кількість ітерації
тесту, яка забезпечує
припустиму ймовірність помилки для застосування у генерації ключів RSA
з встановленим рівнем стійкості захисту. Припустиму
ймовірність помилки пропонується прив'язувати до рівня
стійкості захисту як \(2^{-s}\), де \(s\) -- рівень
стійкості захисту для \(n\)-бітного модуля в схемі RSA
\cite[][\pagename 73]{2023}.
Отже, за такої моделі прийнятність малої ненульової
ймовірності помилки тесту простоти обґрунтовується цільовим
рівнем стійкості захисту.

Можна також припустити, що, за малої ймовірності помилки тесту,
якщо складене число все ж таки не було виявлено, розкласти
його на множники буде також складно.

Як жартівливо зазначає Д. Кнут у \cite[][\pagename 395]{Knuth2021},
для 25 ітерацій, що дає меншу за \(\frac{1}{4}^25\)
ймовірність не виявити складене число,
більш ймовірна апаратна помилка, яка вплине на коректність обчислення.
Хоча, при виявлення складеного числа алгоритм визначає
\enquote{свідка} складеності,
і цей результат можна однозначно підтвердити або спростувати.

\begin{quote}\small\it
  Які вимоги, крім достатньої довжини, висуваються до
  простих чисел, що використовуються при формуванні RSA-ключів?
\end{quote}

Для математичної коректності схема RSA вимагає, щоб для згенерованих
ключів \(p-1\), \(q-1\) були взаємно прості із експонентою.

У \cite{2023} вимагається виконання наступних співвідношень
для згенерованих простих чисел \(p\), \(q\) з бітовою довжиною
\(\mathit{nBits}\) публічного модуля \(n = pq\):
\begin{enumerate}[label=\alph*.]
  \item \( 2^{(\mathit{nBits}-1)/2} < p < 2^{\mathit{nBits}/2}\)
  \item \( 2^{(\mathit{nBits}-1)/2} < q < 2^{\mathit{nBits}/2}\)
  \item \(|p - q| > 2^{\mathit{nBits}/2 - 100}\)
  \item \(\operatorname{GCD}(e, \operatorname{LCM}(p-1,q-1)) = 1\)
  \item \(2^{\mathit{nBits}/2} < d < \operatorname{LCM}(p-1,q-1)\)
    для \(d = e ^{-1} \mod (\operatorname{LCM}(p-1, q-1))\)
\end{enumerate}

У \cite[][\pagename 405]{Knuth2021} для RSA
наводяться наступні вимоги: неподільність \(p-1\), \(q-1\) на 3;
\(p-1\) та \(q-1\)повинні мати принаймні один дуже великий дільник;
відношення \(\frac{p}{q}\) не повинно бути близьким до простого дробу.

Вимоги до згенерованих чисел обумовлені не тільки
математичною коректністю, а й міркуваннями безпеки.

У таблиці \ref{table:primality-test-comparison} порівняно
тести простоти Ферма, Міллера-Рабіна, Соловея-Штрассена.

\bgroup
\renewcommand{\arraystretch}{1.2}
% tex-fmt: off
\begin{longtable}{l||l|l|l}
  \caption{Порівняння ймовірнісних тестів простоти}
  \label{table:primality-test-comparison}
  \\
  Тест                    & \parbox{23ex}{Ймовірність \enquote{помилки}\par для \(k\) ітерацій}  & Принципові недоліки                                   & Оцінка часової складності
  \bigskip\endhead
  Ферма                   & \(2^{-k}\)                                        & \parbox{20ex}{Псевдопрості Ферма,\par числа Кармайкла}                   & \(\mathnormal{O}(k \log^2 n \log\log n)\) \\
  Соловея-Штрассена       & \(2^{-k}\)                                        & --                                                                       & \(\mathnormal{O}(k \log^3 n)\)              \\
  Міллера-Рабіна          & \(4^{-k}\)                                        & --                                                                       & \(\mathnormal{O}(kn^2\log n \log\log n)\) \\
\end{longtable}
% tex-fmt: on
\egroup

Таким чином, серед розглянутих алгоритмів
ймовірнісний тест простоти Міллера-Рабіна є оптимальним для
практичного застосування через нижчу ймовірність похибки та
поліноміальну складність ітерації.

\chapter{Практичне дослідження генерації випадкових значень}

Розглянемо генерацію псевдовипадкових чисел стандартними
засобами мови C.
Стандарт C23 гарантує реалізацію функції для генерації
псевдовипадкових чисел \texttt{rand()} для
некриптографічних застосувань без гарантій якості.

\section{Використання звичайного генератора псевдовипадкових чисел}

\inputminted{c}{3.1.c}
\begin{minted}{shell}
clang -xc -std=c23 -o 3.1 3.1.c
./3.1
sleep 5
./3.1
\end{minted}
\begin{minted}{output}
-- 57:35
Dice roll: 1968078301
Dice roll: 287724083
Dice roll: 410622274
Dice roll: 558519326
Dice roll: 460165363
Dice roll: 773440537
Dice roll: 901520025
Dice roll: 1404622363
Dice roll: 417397028
Dice roll: 665131385
-- 57:40
Dice roll: 1968078301
Dice roll: 287724083
Dice roll: 410622274
Dice roll: 558519326
Dice roll: 460165363
Dice roll: 773440537
Dice roll: 901520025
Dice roll: 1404622363
Dice roll: 417397028
Dice roll: 665131385
\end{minted}

Отже, стандартний інтерфейс генерації псевдовипадкових
чисел для мови C використовує детерміністичний алгоритм, єдиним
джерелом ентропії якого є початкове значення, задане за
допомогою виклику \texttt{srand()}.

Така відтворюваність дозволила б значно зменшити простір
перебору ключів.
Наприклад, якби в якості джерела початкового значення
використовувався б поточний
системний час, то порушних міг би цим скористатися для
оцінки ймовірних початкових значень і навіть уникнути їх
повного перебору. Також початкове значення і внутрішній
стан описуються числом типу \texttt{int}, тобто на
більшості платформа мають не більше 32 бітів.
Більше того, немає жодних гарантій щодо статистичних
якостей згенерованої псевдовипадкової послідовності.

\section{Використання криптографічно стійкого генератора}

Реалізація стандартної бібліотеки C \emph{GNU C Library}
з версії 2.36 надає сімейство функцій \texttt{arc4random*},
які \enquote{надають криптографічно стійкі псевдовипадкові числа},
використовуючи ентропію, надану ядром операційної системи.

\inputminted{c}{3.2.c}
\begin{minted}{shell}
clang -xc -std=c23 -o 3.2 -D_DEFAULT_SOURCE 3.2.c
./3.2
sleep 5
./3.2
\end{minted}
\begin{minted}{output}
0075e28e9656e331687bde4f51406da3c5fc14ddea95813308fabd60d461a213
87fc0cacbbaaa21bca01cd052e01aec71168601ed29ac6a20d06d84bf2f004e6
\end{minted}

Отже, ці розглянуті функції дійсно використовують додаткове
джерело ентропії.
На відміну від попереднього прикладу, за тиз самих умов
генеруються різні значення.

Як задокументовано в \emph{The GNU C Library Reference
Manual, for version 2.44}:
\blockquote{
  Бібліотека GNU C використовує ентропію з ядра,
  отриману через getrandom або /dev/urandom для початкових значень
}.

\chapter{Дослідження генерації асиметричних ключів}

Для дослідження генерації асиметричних ключів обрано
криптографічну бібліотеку \texttt{OpenSSL},
яка написана мовою C та має сертифікати відповідності FIPS-140.

\inputminted{c}{4.1.c}

\begin{minted}{shell}
clang -xc -std=c23 -o 4.1 -lcrypto 4.1.c
\end{minted}

\begin{minted}{shell}
./4.1 -kT 2048
\end{minted}
\begin{minted}{output}
Generating key of length: 2048
Generation took 87685 clocks
Private parameters:
  Private-Key: (2048 bit, 2 primes)
  modulus:
      00:ab:77:4e:34:44:eb:26:9c:4d:5d:39:ec:6c:d7:
      4b:6a:f1:6f:73:23:83:a0:2b:1f:18:3f:86:e4:55:
      2d:08:0b:7e:15:78:52:ec:4f:44:60:1b:7f:a6:45:
      82:77:d4:a6:e1:7b:fe:2c:04:fd:f8:f9:6e:2c:50:
      b0:23:2e:bd:c6:e3:90:43:91:7a:23:86:2f:68:92:
      3c:fa:b2:24:cd:87:86:52:20:c4:0f:e5:28:1f:dd:
      07:05:42:db:18:2f:3e:fc:cd:83:0e:76:8a:15:67:
      f9:54:8e:3c:3d:ad:03:99:cf:8a:3d:f8:66:d1:04:
      e6:9e:6c:d6:bc:b9:4a:83:4c:df:ba:d0:86:ac:28:
      ba:9c:c5:d6:c5:d5:5a:cc:55:de:15:c0:3a:99:c9:
      db:2a:34:be:f5:7a:de:a7:d3:bf:fa:e2:88:ec:a4:
      c2:20:26:b3:b2:ca:0f:be:1e:74:7a:b9:b6:45:f2:
      89:b2:d8:31:aa:69:08:f7:9e:35:99:4b:ab:c1:a2:
      c7:79:21:05:03:13:f2:d8:e5:95:a9:b6:a7:8d:3b:
      4a:e4:0b:33:14:a0:70:f2:c0:16:29:ac:dc:64:c5:
      25:67:71:46:ff:cb:49:14:d7:23:61:a5:ce:96:15:
      5c:64:f5:0c:6e:c9:64:8d:55:25:b1:85:fe:70:9e:
      8e:e9
  publicExponent: 65537 (0x10001)
  privateExponent:
      05:9b:44:f6:da:c4:0a:cf:da:f2:89:85:9f:9f:73:
      be:8c:25:09:44:a7:e9:0b:1c:b6:ac:29:17:ae:03:
      4d:6e:89:f7:f8:c0:1f:c6:27:58:e5:8a:16:13:9f:
      91:62:e6:28:cf:80:cb:fe:ad:64:73:99:0e:54:17:
      e3:19:d0:db:59:ae:ea:9b:7a:3d:02:86:1d:51:85:
      32:b8:1a:44:f6:8e:5f:b4:18:5e:ae:af:05:49:4c:
      5e:5e:5e:f6:34:2e:14:2b:75:4f:46:db:2e:66:9a:
      6f:2a:a7:d7:cd:a0:8b:d0:75:24:39:cb:2b:a7:11:
      12:2c:db:e1:d7:ab:2c:8e:72:9b:6e:22:21:d7:f5:
      f0:43:24:50:64:cf:49:e0:eb:7a:9e:5e:e4:3e:63:
      23:2e:a9:25:3c:83:08:25:69:9f:63:2a:9c:1a:e0:
      2d:18:65:f8:08:be:34:8f:7e:8b:81:75:8f:01:0f:
      ab:c8:84:ff:a0:88:f6:ce:ba:36:37:35:16:66:98:
      a1:66:0e:17:8b:e4:91:53:1d:a5:0b:27:8b:9b:00:
      af:16:09:7f:d7:c0:81:55:95:68:0f:cb:d2:9e:66:
      39:c1:a8:62:06:9b:d3:d4:a7:85:f8:29:4f:50:60:
      23:96:98:cc:1f:49:e6:d4:69:e0:5b:bd:a3:e1:35:
      4b
  prime1:
      00:ef:ac:3e:a9:a4:c2:45:d9:2b:02:0a:52:52:d4:
      9e:21:8f:19:e0:cc:fc:16:71:1d:80:9d:ea:1c:e4:
      a7:b6:55:21:67:6a:51:b4:dd:3a:97:4a:73:a4:71:
      79:5a:4a:43:53:1e:16:72:9c:32:64:27:9b:d9:ad:
      f8:19:7d:ce:4e:7d:91:46:9c:de:96:91:d4:d6:35:
      9d:1d:c6:39:32:2e:92:8c:1f:ff:2f:28:6b:44:e5:
      70:39:81:b0:d0:16:7f:82:82:b9:85:2c:df:72:2b:
      3d:25:56:51:f1:ce:f1:e2:e7:52:b1:bd:6d:ea:8f:
      bc:cf:45:f5:7c:33:ba:9e:3f
  prime2:
      00:b7:25:92:e2:21:fe:6a:b3:d3:6f:9d:17:e8:dc:
      4c:b6:a4:cb:f1:e9:83:79:d2:92:9f:77:83:eb:14:
      02:18:f6:7a:4d:05:79:5a:b3:9d:b2:ea:af:02:56:
      70:e8:5f:69:52:7b:9d:1f:9f:f6:06:74:39:0c:d1:
      31:ad:77:64:2b:3e:c7:2e:4f:a9:e9:99:10:d7:77:
      46:3f:fb:f4:da:36:df:2f:b2:5d:1c:02:53:8b:cd:
      e3:c4:0a:87:c3:6e:46:1b:de:44:ca:bc:e7:ac:ff:
      ee:7b:a5:6a:18:e2:20:96:d2:d4:f2:f8:bd:28:dc:
      37:05:70:4e:1f:64:02:58:d7
  exponent1:
      2c:07:f0:db:29:b1:a7:33:66:7b:51:7a:1e:4d:81:
      5f:98:f2:a0:22:08:8d:f6:e9:fc:9f:06:4b:98:96:
      57:20:fc:aa:0b:16:47:ab:fc:ff:9e:65:f0:68:08:
      6b:56:09:58:53:c6:bb:f7:86:2a:70:3c:69:99:54:
      63:1a:81:b4:59:cd:3e:b3:c3:10:d6:69:cb:90:bd:
      f4:a2:5e:0f:a2:0e:ba:3f:86:9a:1b:4a:d9:b5:6d:
      10:a3:8a:e0:2f:c1:ed:30:26:f5:18:71:54:80:2d:
      13:96:f0:27:04:65:db:d1:b3:97:83:bd:95:c2:22:
      65:d1:ec:6f:8c:f1:53:f1
  exponent2:
      00:8f:c6:37:60:be:3e:f3:94:e7:cf:1f:0b:be:09:
      f3:4f:0e:5f:ec:30:dc:bd:ad:c8:dc:d6:fe:00:75:
      5c:90:2c:2e:26:78:bd:a3:ae:8d:85:ff:26:14:37:
      b7:bb:33:eb:82:8c:f7:f3:f3:a7:37:b9:d5:6f:72:
      6a:05:3b:4c:af:86:40:e6:e2:33:08:e6:d8:89:cd:
      60:a0:aa:32:b5:7b:8e:09:eb:83:9c:91:56:f4:d1:
      54:94:79:8f:ca:a3:ae:f6:01:af:0f:58:ec:a6:aa:
      e3:94:bd:6e:64:ca:5f:bb:07:17:95:06:1f:b9:0b:
      e7:45:a5:bc:2e:78:5c:3f:51
  coefficient:
      7c:92:fa:b7:51:59:65:e2:61:77:87:4b:8a:c1:70:
      7c:a2:73:ed:ab:6b:64:c2:f6:71:d0:26:5a:a8:4e:
      4e:f4:3b:d9:3c:b7:ad:4a:93:b2:34:f5:9a:aa:30:
      8b:63:d8:59:4b:bc:cb:5f:cf:6f:c6:c1:76:32:40:
      c2:54:1c:04:a0:ca:b9:7a:71:de:d0:f1:8d:f5:99:
      30:ec:f5:17:28:f0:b4:0a:17:bd:cd:10:a8:31:3f:
      dc:ea:fd:32:91:60:92:b7:06:23:fa:bf:35:67:e3:
      1f:32:f2:e8:c1:00:8d:51:64:ba:95:99:04:ee:6a:
      07:50:69:4b:8b:5b:eb:af
\end{minted}

\inputminted{bash}{4.1.sh}
\begin{minted}{shell}
./4.1.sh
\end{minted}
\begin{minted}{output}
Timing 100 generations of 2048-bit RSA keys

real    0m3,915s
user    0m3,749s
sys     0m0,150s
Timing 100 generations of 4096-bit RSA keys

real    0m32,142s
user    0m31,898s
sys     0m0,151s
\end{minted}

Отже, приблизний середній час генерації ключа з модулем довжиною
2048 біт -- 39 мс,
ключа з модулем довжиною 4096 біт -- 321 мс.

Час генерації не є сталим навіть для тої самої довжини ключа, оскільки
процес генерації є випадковим, а також залежить від
накопичення достатнього рівня ентропії.

Для ключів із більшим модулем простір пошуку простих чисел
для формування модуля також збільшується. Крім того,
необхідна більша кількість ітерацій тестів на простоту для
гарантування прийнятного рівня впевненості, що знайдені числа прості.
Інші перевірки також займають більше часу, оскільки
операції реалізовані через арифметику великих чисел
залежать від довжини бітового представлення числа.

Як вже було пояснено раніше, надійність CSPRNG та
тестів на простоту при генерації ключів відіграють ключову
роль для коректності і безпеки застосування схеми RSA.

У розділі \ref{chapter:rsa-keygen} наведені вимоги,
необхідність дотримання яких робить реалізацію коректного і безпечного
алгоритма генерації ключів RSA складною і чутливою до помилок задачею.
Розробка алгоритма генерації ключів з відповідним рівнем
безпекових гарантії
фактично означає розробку цілої бібліотеки криптографічних примітивів.
Зважаючи на це та критичність сфери застосування,
розробники прикладного програмного забезпечення не повинні
намагатися реалізувати такий алгоритм самостійно.

\chapter{Аналіз ситуацій порушення безпеки}

\section{
  \textbf{Ситуація 1.} Передбачувана генерація ключа
}

Розмір простору можливих ключів  в такому випадку зумовлений
кількістю можливих значень \texttt{current\_timestamp}.
Наприклад, якщо \texttt{current\_timestamp} є деяким
представленням POSIX timestamp, що описується 64-бітовим
значенням, то фактичний розмір ключового простору лише \(2^64\).
Крім того, оскільки \texttt{current\_timestamp} відповідає
календарному часу,
то для конкретного застосування розмір простора буде ще меншим.
Якщо \texttt{current\_timestamp} -- це рядок з
представленням з точністю до мілісекунд,
то за певний проміжок часу можлива кількість різних ключів
-- це лише довжина цього проміжку часу в мілісекундах.
Тобто, обране джерело ентропії (лічильник часу) обмежує
рівень стійкості захисту ключа.
Також, лічильник часу не є випадковим, і будь-хто може
точно оцінити можливі значення
\texttt{current\_timestamp} та перебрати можливі згенеровані ключі.

\section{
  \textbf{Ситуація 2.} Однаковий seed на різних пристроях
}

Схоже на попередню ситуацію, в такому разі дослідник
простір ключів фактично дуже малий.
Навіть якщо припустити, що генератор псевдовипадкових чисел
не завжди опиняється в тому самому стані на кожному
пристрої на момент генерації ключів,
дослідник всеодно може оцінити можливі внутрішні стани генератора
псевдовипадкових чисел на момент створення ключа та суттєво
зменшити простір перебору для довільного пристрою.

\section{
  \textbf{Ситуація 3.} Розголошення правильно згенерованого ключа
}

В такому разі система не є надійною, бо значення ключа стає
відомим не лише його власнику, і ключ будь-який потрібно вважати
скомпрометованим.

\section{
  \textbf{Ситуація 4.} Недостатньо випадкова генерація RSA-ключа
}

В такому разі дослідник зможе легко відновити будь-яку пару
ключів, які використовували одне й те саме просто число.
Для цього достатньо було б обрахувати НСД публічних модулів
цих ключів для знаходження одного з секретних простих
множників. Фактично, схема RSA в такому разі є повністю
скомпрометованою.

\section{
  \textbf{Ситуація 5.} Надмірно тривале використання одного ключа
}

Безпека асиметричної криптографії будується на складності
розв'язання односторонніх задач, але цей розв'язок
можливо знайти завдяки повному перебору. Довготривале
використання одного ключа створює загрозу зламу ключа,
в тому числі через можливе виявлення вади даного конкретного ключа.
Наприклад, оскільки тести простоти при генерації RSA ключів
ймовірнісні, то існує ризик, що модуль згенерованого ключа
розкладається більше ніж на два множники, і це посприяє зламу.
Тому регулярна ротація значно знижує такі ризики.

Залежно від використаної криптосистеми, довготривале
застосування того самого ключа створює можливість для
накопиченням криптографічного матеріалу для аналізу.

Безпека деяких схем залежить від недопущення повторного
використання одноразових значень, або може бути безпечно
застосована лише до обмеженого обсягу даних \cite{Barker2020a}.
При порушенні умов безпека системи буде скомпрометована.

У \cite{DiCrescenzo1999} навіть розглядаються випадки, коли
довготривале зберігання ключи ускладнює задачу його повного
знищення на фізичному носії.

\chapter{Життєвий цикл криптографічного ключа}

Розглянемо аналіз загроз безпеки на різних етапах життєвого циклу
криптографічного ключа. Розглядається наступна модель
життєвого циклу ключа:
{\it
  генерація
  $\Longrightarrow$ зберігання
  $\Longrightarrow$ використання
  $\Longrightarrow$ ротація/заміна
  $\Longrightarrow$ відкликання або знищення
}.

\section{Генерація}
\subsection{Загрози}
\begin{enumerate}
  \item
    Низька ентропія, передбачуваність згенерованих даних.
  \item
    Обрані небезпечні параметри криптографічних алгоритмів.
  \item
    Згенерований ключ можу виявитися слабким.
  \item
    Витік інформації про ключ сторонніми каналами.
\end{enumerate}
\subsection{Вимоги безпеки}
\begin{enumerate}
  \item
    Використовувати надійні, перевірені реалізації
    програмних, апаратних криптографічних засобів, зокрема,
    для генерації випадкових бітів.
  \item
    Обирати параметри криптографічних алгоритмів згідно
    з актуальними  практиками безпеки, рекомендаціями
    авторитетних інституцій.
  \item
    Тестувати згенерований ключ на відомі класи безпекових вад.
\end{enumerate}

\section{Зберігання}
\subsection{Загрози}
\begin{enumerate}
  \item
    Викрадення ключів.
  \item
    Випадкове знищення ключів.
\end{enumerate}
\subsection{Вимоги безпеки}
\begin{enumerate}
  \item
    Зберігання ключів на спеціалізованих носіях;
    криптографічний захист конфіденційності і цілісності
    самих ключів із дотриманням  достатнього рівня
    стійкості захисту, наприклад, згідно з
    рекомендаціями \cite{Barker2020a}.
  \item
    Резервне копіювання, архівне зберігання ключового матеріалу із
    збереженням відповідного рівня захисту.
  \item
    Дотримання встановленого криптоперіоду.
  \item
    Належне розмежування доступу до ключового
    матеріалу, зокрема, програмними, апаратними,
    організаційними засобами.
\end{enumerate}

\section{Використання}
\subsection{Загрози}
\begin{enumerate}
  \item
    Некоректне використання, що призводить до компрометації ключа.
  \item
    Накопичення криптографічного матеріалу для даного
    ключа зловмисником для підсилення можливостей аналізу.
\end{enumerate}
\subsection{Вимоги безпеки}
\begin{enumerate}
  \item
    Дотримання вимог і обмежень використання
    криптографічних алгоритмів.
  \item
    Розмежування призначення ключа (для шифрування, для
    підписів \dots)
  \item
    Дотримання встановленого криптоперіоду і своєчасна
    ротація ключів.
  \item
    Аудит використання, доступу до ключового матеріалу.
\end{enumerate}

\section{Ротація/Заміна}
\subsection{Загрози}
\begin{enumerate}
  \item
    Заміна ключа на пов'язаний ключ, з подальшим
    розкриттям історичного ключа
    (key update function у \cite{Barker2020a}).
  \item
    Замінений ключ приймається деякими сторонами як дозволений.
  \item
    Несанкціонована ініціація процедури ротації.
  \item
    Порушення доступності до даних, захищених старим ключом.
\end{enumerate}
\subsection{Вимоги безпеки}
\begin{enumerate}
  \item
    Дотримання вимог генерації ключа.
  \item
    Незалежність нового ключа.
  \item
    Своєчасне, безпечне інформування зацікавлених
    сторін про ротацію ключа.
  \item
    Контролювати реалізацію дозволу на ініціалізацію ротації ключа.
  \item Перехідний період і належна зміна активного стану
    старого ключа, збереження метаданих тта дозвіл,
    зокрема, перевіряти підписані старим ключом дані.
\end{enumerate}

\section{Відкликання або знищення}
\subsection{Загрози}
\begin{enumerate}
  \item
    Замінений ключ приймається деякими сторонами як дозволений;
    затримка фактичного припинення повноважень ключа.
  \item
    Ключ може бути відновлений, зокрема з резервних
    копій, архівних матеріалів, з фізичного обладнання.
  \item
    Неможливість проаналізувати випадки використання
    знищеного ключа, відсутність належної можливості для аудиту.
\end{enumerate}

\subsection{Вимоги безпеки}
\begin{enumerate}
  \item
    Своєчасне, безпечне інформування зацікавлених
    сторін про відкликання ключа.
  \item
    Гарантоване знищення всіх записів ключа.
  \item
    Збереження метаданих ключа в цілях аудиту.
\end{enumerate}

\chapter{Висновки}

Безпека криптографічних систем з використанням симетричних
та асиметричних ключів залежить, зокрема, від умов генерації ключового
матеріалу, дотримання вимог безпеки життєвого циклу ключа.

Критично важливим примітивом для створення криптографічних
ключів є генератор випадкових бітів. Рівень стійкості
захисту ключа та системи в цілому залежить від рівня
захисту використаного генератора.

Безпека будь-якого генератора випадкових бітів, генератора
псевдовипадкових чисел залежить обмежується доступною ентропією.

У криптосистемах, безпека яких ґрунтується на складності
задачі факторизації, застосовуються ймовірнісні тести простоти чисел,
що зумовлено практичними обмеженнями на використання детерміністичних
алгоритмів. Тому модель безпеки повинна передбачати
ймовірнісний характер гарантій властивостей ключів в таких
криптосистемах. Це, зокрема, впливає на застосовану модель
життєвого цикла ключа, обмеження його криптоперіоду для
зниження ризику компрометації ключа.

В життєвому цикл криптографічного ключа можна виділити
наступні етапи: генерація, зберігання, використання,
ротація або заміна, відкликання або знищення.
Для кожного етапу було розглянуто можливі загрози та вимоги
безпека для протидії цим загрозам.

Безпечна реалізація і експлуатація гібридних
криптографічних систем повинна ґрунтуватися на
використанні надійних програмних і апаратних засобів;
строгому дотриманні вимог криптографічних протоколів;
інтеграції заходів з мінімізації ризиків в організаційні політики;
дотриманні актуальних, визнаних практик безпеки,
рекомендованих, зокрема, міжнародними і відкритими стандартами.

\printbibliography[
  heading=bibintoc,
  title={Перелік використаних джерел}
]

\end{document}
